% Section 20: Parallel transport} | Part: 4 | Depends on: see docs/REWRITE_PLAN.md | Status: todo | Original lines: 2998-3020
%!TEX root = ../../../main.tex
\subsection{Parallel transport}
\begin{formal}[{parallel transport along a curve}]
For a smooth curve $\gamma(t)$ and a vector field $Y$ along $\gamma$ (a smooth assignment $t\mapsto Y_{\gamma(t)}\in T_{\gamma(t)}M$), define $\nabla_{\dot\gamma}Y:=\nabla_{\hat X}\hat Y$ for any extensions $\hat X,\hat Y$ of $\dot\gamma, Y$ to fields on $M$; the result does not depend on the extensions. In coordinates $\gamma=(x^i(t))$:
\[
(\nabla_{\dot\gamma}Y)^k=\frac{\dd Y^k}{\dd t}+\Gamma^k{}_{ij}(\gamma(t))\,\dot x^i\,Y^j .
\]
\textbf{Theorem.} Given $v_0\in T_{\gamma(t_0)}M$ there is a unique field $Y$ along $\gamma$ with $Y_{\gamma(t_0)}=v_0$ and $\nabla_{\dot\gamma}Y=0$. (In coordinates this is the linear ODE
\begin{equation}\label{eq:ptode}
\frac{\dd Y^k}{\dd t}+\Gamma^k{}_{ij}(\gamma(t))\,\dot x^i\,Y^j=0,
\end{equation}
so existence and uniqueness follow from ODE theory; a general curve is split into finitely many chart pieces.) The \emph{parallel transport} is the linear isomorphism $\Pi^\gamma_{\gamma(t_0),\gamma(t_1)}:T_{\gamma(t_0)}M\to T_{\gamma(t_1)}M$, $v_0\mapsto Y_{\gamma(t_1)}$.
\end{formal}
\begin{intuition}[{carrying an arrow without turning it}]
Like carrying a spear while walking and trying not to rotate it ``relative to the ground''. On a curved surface this is ambiguous globally, but infinitesimally it is perfectly well defined by \eqref{eq:ptode}. In a flat space with $\nabla^0$ and Cartesian coordinates ($\Gamma=0$) parallel transport leaves components unchanged ($\Pi(v_0)=v_0$ via $T_p\R^m\cong\R^m$).
\end{intuition}

\begin{example}[{parallel transport on the sphere: path-dependence}]
For the unit sphere (metric $\dd\theta^2+\sin^2\theta\,\dd\phi^2$; Christoffel symbols computed in Section~\ref{sec:lc}): $\Gamma^\theta{}_{\phi\phi}=-\sin\theta\cos\theta$, $\Gamma^\phi{}_{\theta\phi}=\cot\theta$. 
\emph{Along the equator} $\theta=\pi/2$, $\dot\theta=0$, $\dot\phi=1$: both symbols vanish at $\theta=\pi/2$, so \eqref{eq:ptode} gives $\dot Y^\theta=\dot Y^\phi=0$: the components $(Y^\theta,Y^\phi)$ stay constant.
\emph{Along a meridian} ($\phi=\mathrm{const}$, $\dot\theta=1$): $\dot Y^\theta=0$ and $\dot Y^\phi+\cot\theta\,Y^\phi=0$, so $Y^\theta$ is constant, $Y^\phi\propto 1/\sin\theta$ (hence its length $\sin\theta\,|Y^\phi|$ is constant).
\emph{A triangle.} Go from the north pole down the meridian $\phi=0$ to the equator, along the equator to $\phi=\pi/2$, and back up to the pole. Start at the pole with a vector pointing along the first meridian's direction ``away from the pole'' ($\theta$-direction at $\phi=0$). Transported down the first meridian it remains tangent to it; along the equator it stays pointing towards the south (constant $Y^\theta$), perpendicular to the path; going back up the second meridian it remains the same way, now pointing \emph{along} the second meridian, away from the pole. At the end the vector has rotated by $\pi/2$ with respect to the initial one. The rotation angle equals the enclosed area times the curvature ($K=1$): the octant has area $4\pi/8=\pi/2$. \textbf{Parallel transport depends on the path precisely because the space is curved} (Section~\ref{sec:curv}).
\end{example}
